\begin{lemma}
In a $3$-uniform $3$-simple local extremal setup at scale $D$, with
$\mathcal F$, $f$, $P(f)$, $T(f)$, and
$b_{L(\mathcal F),3D}(f)$ as in
\noderef{ThreeUniformThreeSimpleLocalSetup}, the independent-pair count
satisfies
\[
P(f)\le 1.031D^2-9D
\]
for all sufficiently large $D$.
\end{lemma}
\begin{proof}
Let $G$ be the subgraph of $L(\mathcal F)$ induced by
$N_{L(\mathcal F)}(f)$, and let $\overline G$ be its complement.  The
independent pairs and independent triples in the line-graph neighborhood are
the edges and triangles of $\overline G$, by the definitions
\noderef{IndependentPairCount}, \noderef{IndependentTripleCount}, and
\noderef{LineGraphOfHypergraph}.  The graph $\overline G$ is tripartite:
the saturated edge structure separates the neighbors of $f$ according to the
unique vertex of $f$ that they contain, and two neighbors in the same class
intersect at that vertex, so they are adjacent in $G$ and nonadjacent in
$\overline G$.  Hence \noderef{TriangleCountPartite} gives
\[
T(f)\le \left(\frac{P(f)}3\right)^{3/2}.
\]

Assume for contradiction that
$P(f)>1.031D^2-9D$.  Since
\noderef{ExtremalBParameterSaturationThreeSimple} gives
$\deg_{L(\mathcal F)}(f)=3(D-1)$ in the extremal setup, the formula
\noderef{LocalBParameter} and the triangle bound just proved give
\[
b_{L(\mathcal F),3D}(f)
\le
\frac{3(D-1)}{3D}
-\frac{P(f)}{(3D)^2}
+\frac{(P(f)/3)^{3/2}}{(3D)^3}.
\]
Put
\[
p=\frac{P(f)}{D^2}.
\]
Because $P(f)$ is an edge count of $\overline G$, it is nonnegative.  Also
$\overline G$ has the same vertex set as the neighborhood of $f$, and
\noderef{ThreeUniformThreeSimpleLocalSetup} together with
\noderef{ExtremalBParameterSaturationThreeSimple} says that this neighborhood
has $3(D-1)$ vertices.  Hence
\[
0\le P(f)\le \binom{3D-3}{2},
\qquad\text{so}\qquad
0\le p\le \frac{(3D-3)(3D-4)}{2D^2}\le \frac92
\]
for every positive $D$.  In terms of $p$, the displayed upper bound is
\[
b_{L(\mathcal F),3D}(f)
\le
-\frac1D+F(p),
\qquad
F(x)=1-\frac{x}{9}+\frac{x^{3/2}}{81\sqrt3}.
\]
For $0\le x\le 9/2$,
\[
F'(x)=-\frac19+\frac{\sqrt{x}}{54\sqrt3}<0,
\]
because $\sqrt{x}\le\sqrt{9/2}<6\sqrt3$.  Thus $F$ is decreasing on the
whole admissible interval.  Moreover $F'(x)\ge -1/9$ there, since the second
term in $F'(x)$ is nonnegative.

The contradiction hypothesis says
\[
p>1.031-\frac9D.
\]
For all sufficiently large $D$ the left endpoint
$a_D=1.031-9/D$ is positive, and $a_D\le 1.031$.  By monotonicity,
$F(p)<F(a_D)$.  By the lower derivative bound, or equivalently by the mean
value theorem applied on $[a_D,1.031]$,
\[
F(a_D)-F(1.031)\le \frac{1.031-a_D}{9}=\frac1D.
\]
Therefore
\[
b_{L(\mathcal F),3D}(f)
\le -\frac1D+F(p)
< -\frac1D+F(a_D)
\le F(1.031).
\]
It remains only to record the numerical margin in a checkable form.  Since
$1.031=1031/1000$,
\[
F(1.031)
=1-\frac{1031}{9000}
  +\frac{1031\sqrt{1031}}{810000\sqrt{30}}.
\]
The last term is at most $6719/900000$: both sides are nonnegative, and after
squaring this is exactly the rational inequality
\[
\frac{1031^3}{810000^2\cdot 30}
<
\frac{6719^2}{900000^2},
\]
whose difference is
$11097613/196830000000000>0$.  Hence
\[
F(1.031)
<
1-\frac{1031}{9000}+\frac{6719}{900000}
=0.89291.
\]
Thus the assumption $P(f)>1.031D^2-9D$ forces
\[
b_{L(\mathcal F),3D}(f)<0.89291.
\]

On the other hand, apply
\noderef{ThreeUniformThreeSimpleExtremalExample} with
$\eta=1/300000$.  It gives, for all sufficiently large $D$, an admissible
$3$-uniform $3$-simple hypergraph-edge pair $(\mathcal H,e)$ with
\[
b_{L(\mathcal H),3D}(e)\ge
1-\frac19+\frac1{3^5}-\frac1{300000}.
\]
The right side is
\[
1-\frac19+\frac1{243}-\frac1{300000}
=0.893000781\ldots>0.893.
\]
Since \noderef{ThreeUniformThreeSimpleLocalSetup} includes the choice of
$(\mathcal F,f)$ maximizing $b_{L(\mathcal H),3D}(e)$ over all admissible
$3$-uniform $3$-simple hypergraphs of maximum degree at most $D$, we get
\[
b_{L(\mathcal F),3D}(f)\ge b_{L(\mathcal H),3D}(e)>0.893.
\]
This contradicts the upper bound
$b_{L(\mathcal F),3D}(f)<0.89291$.  Therefore the assumed strict inequality
for $P(f)$ is impossible, and the cutoff follows.
\end{proof}
